\documentclass[10pt]{beamer}
\usetheme{metropolis}
\usepackage{graphicx}
\usepackage{listings}
\usepackage{booktabs}
\usepackage{xcolor}
\usepackage{hyperref}

\definecolor{codebg}{HTML}{F5F5F5}
\definecolor{codegreen}{HTML}{2E7D32}
\definecolor{codegray}{HTML}{757575}
\definecolor{codeblue}{HTML}{1565C0}

\lstdefinestyle{rstyle}{
  language=R, backgroundcolor=\color{codebg}, basicstyle=\ttfamily\scriptsize,
  keywordstyle=\color{codeblue}\bfseries, stringstyle=\color{codegreen},
  commentstyle=\color{codegray}\itshape, breaklines=true, frame=single,
  rulecolor=\color{codegray!30}, showstringspaces=false, tabsize=2,
  morekeywords={library,ggplot,aes,geom_point,geom_line,geom_col,geom_bar,
    geom_histogram,geom_boxplot,geom_smooth,geom_density,geom_jitter,
    geom_text,geom_label,geom_segment,geom_hex,geom_violin,geom_area,
    geom_tile,geom_rug,geom_text_repel,
    scale_color_manual,scale_fill_viridis_c,scale_color_brewer,
    scale_size,scale_shape_manual,scale_x_log10,
    facet_wrap,coord_flip,theme_minimal,theme,element_text,labs,
    position_jitter,after_stat,stat_summary}
}
\lstdefinestyle{pythonstyle}{
  language=Python, backgroundcolor=\color{codebg}, basicstyle=\ttfamily\scriptsize,
  keywordstyle=\color{codeblue}\bfseries, stringstyle=\color{codegreen},
  commentstyle=\color{codegray}\itshape, breaklines=true, frame=single,
  rulecolor=\color{codegray!30}, showstringspaces=false, tabsize=4,
  morekeywords={import,from,as,ggplot,aes,geom_point,geom_line,geom_col,
    geom_smooth,geom_histogram,geom_boxplot,geom_bar,geom_density,
    geom_text,geom_label,geom_segment,geom_jitter,
    scale_color_brewer,scale_color_cmap,scale_size,
    facet_wrap,theme_minimal,labs,coord_flip}
}

\title{ggplot2 Deep Dive: Aesthetics \& Geometries}
\subtitle{Workshop 2 --- Module 3}
\author{Prof.Asoc.\ Endri Raco}
\institute{Data Visualization for Data Scientists \\ UNYT Tirana}
\date{2025/26}

\begin{document}

\begin{frame}
  \titlepage
\end{frame}

\begin{frame}{Module Roadmap}
  \begin{enumerate}
    \item The five core aesthetic channels: x/y, color, size, shape, alpha
    \item Continuous vs discrete colour mapping
    \item Layering multiple geometries with \texttt{+}
    \item Overplotting: diagnosis and three fixes
    \item Special geometries: text, label, segment, lollipop, dumbbell
  \end{enumerate}
  \begin{exampleblock}{Learning Outcome}
    You will be able to encode up to five variables simultaneously using
    aesthetic channels, layer multiple geoms on one axes, and solve
    overplotting with alpha, jitter, or hexbin.
  \end{exampleblock}
\end{frame}

\section{The Five Core Aesthetics}

\begin{frame}{Aesthetic Channels}
  \begin{center}
    \includegraphics[width=0.82\textwidth]{figures_w02m03/core_aesthetics}
  \end{center}
  \begin{alertblock}{Pro-Tip}
    Position (x, y) is the most precise channel (Cleveland rank 1).
    Colour hue is next for categorical data. Size works for ordered
    quantitative data but not for precise comparisons. Shape tops out
    at 6 distinguishable glyphs. Alpha is a last resort.
  \end{alertblock}
\end{frame}

\begin{frame}{The shape Palette: Maximum 6 Levels}
  \begin{center}
    \includegraphics[width=0.85\textwidth]{figures_w02m03/shape_palette}
  \end{center}
  \vspace{4pt}
  ggplot2's default shape scale uses filled shapes for the first 6 levels,
  then switches to unfilled --- which become nearly indistinguishable.
  If your categorical variable has $>6$ levels, use colour or faceting
  instead.
\end{frame}

\begin{frame}{Continuous vs Discrete Colour}
  \begin{center}
    \includegraphics[width=0.92\textwidth]{figures_w02m03/continuous_vs_discrete}
  \end{center}
  \begin{exampleblock}{Data Insight}
    Continuous variables map to a colour gradient
    (\texttt{scale\_color\_viridis\_c}). Discrete variables map to a
    qualitative palette (\texttt{scale\_color\_brewer}). Mismatching
    produces misleading perceptual ordering.
  \end{exampleblock}
\end{frame}

\begin{frame}[fragile]{Multi-Aesthetic Encoding in R}
  \begin{columns}[T]
    \begin{column}{0.52\textwidth}
\begin{lstlisting}[style=rstyle]
# Five variables in one chart
ggplot(df, aes(
    x     = gdp,        # position
    y     = life_exp,    # position
    color = continent,   # hue
    size  = population,  # area
    shape = income_group # glyph
  )) +
  geom_point(alpha = 0.6) +
  scale_color_brewer(
    palette = "Set2") +
  scale_size(range = c(1, 12)) +
  scale_shape_manual(
    values = c(16, 17, 15)) +
  theme_minimal() +
  labs(title = "Gapminder: 5 Vars")
\end{lstlisting}
    \end{column}
    \begin{column}{0.45\textwidth}
      \includegraphics[width=\textwidth]{figures_w02m03/r_multi_aes}

      \vspace{4pt}
      {\small Five variables encoded:
      x = GDP, y = life expectancy,
      colour = continent, size = pop,
      shape = income group. Each
      \texttt{aes()} creates its own
      legend.}
    \end{column}
  \end{columns}
\end{frame}

\begin{frame}[fragile]{Multi-Aesthetic Encoding in Python}
  \begin{columns}[T]
    \begin{column}{0.52\textwidth}
\begin{lstlisting}[style=pythonstyle]
# plotnine: identical syntax
(ggplot(df, aes(
    x="gdp", y="life_exp",
    color="continent",
    size="population",
    shape="income_group"))
 + geom_point(alpha=0.6)
 + scale_color_brewer(
     type="qual", palette="Set2")
 + scale_size(range=(1, 12))
 + theme_minimal())

# matplotlib OO alternative:
for grp in groups:
    mask = df["continent"] == grp
    ax.scatter(
        df.loc[mask, "gdp"],
        df.loc[mask, "life_exp"],
        s=df.loc[mask, "pop"] / 1e5,
        c=palette[grp],
        marker=markers[grp],
        alpha=0.6, label=grp)
\end{lstlisting}
    \end{column}
    \begin{column}{0.45\textwidth}
      \includegraphics[width=\textwidth]{figures_w02m03/py_multi_aes}

      \vspace{4pt}
      {\small plotnine handles multi-
      aesthetic legends automatically.
      In matplotlib, you must loop over
      groups and manage legends manually.}
    \end{column}
  \end{columns}
\end{frame}

\begin{frame}{Bubble Chart: The size Aesthetic}
  \begin{center}
    \includegraphics[width=0.72\textwidth]{figures_w02m03/size_aesthetic}
  \end{center}
  \begin{alertblock}{Pro-Tip}
    \texttt{scale\_size(range = c(1, 12))} controls the min/max point
    size. ggplot2 maps to area (perceptually correct). In matplotlib,
    \texttt{s} is in points\textsuperscript{2} --- already area-based.
  \end{alertblock}
\end{frame}

\section{Layering Multiple Geometries}

\begin{frame}{The Power of \texttt{+}: Three Geom Combinations}
  \begin{center}
    \includegraphics[width=0.95\textwidth]{figures_w02m03/geom_combinations}
  \end{center}
\end{frame}

\begin{frame}[fragile]{Geom Layering in Code}
  \begin{columns}[T]
    \begin{column}{0.48\textwidth}
      \textbf{R: point + smooth}
\begin{lstlisting}[style=rstyle]
ggplot(df, aes(x, y)) +
  geom_point(alpha = 0.4) +
  geom_smooth(method = "lm",
    color = "#E53935")
# Two layers: raw data + trend
\end{lstlisting}
      \vspace{4pt}
      \textbf{R: boxplot + jitter}
\begin{lstlisting}[style=rstyle]
ggplot(df, aes(x = group,
               y = value)) +
  geom_boxplot(outlier.shape = NA,
    fill = "#E3F2FD") +
  geom_jitter(width = 0.15,
    alpha = 0.3, size = 1,
    color = "#E53935")
# Show summary AND raw data
\end{lstlisting}
    \end{column}
    \begin{column}{0.48\textwidth}
      \textbf{R: col + errorbar}
\begin{lstlisting}[style=rstyle]
df_summary <- df |>
  group_by(cat) |>
  summarise(
    m = mean(val),
    se = sd(val)/sqrt(n()))

ggplot(df_summary,
  aes(x = cat, y = m)) +
  geom_col(fill = "#1565C0",
    alpha = 0.6, width = 0.5) +
  geom_errorbar(
    aes(ymin = m - se,
        ymax = m + se),
    width = 0.2,
    linewidth = 0.8) +
  coord_flip()
\end{lstlisting}
    \end{column}
  \end{columns}
\end{frame}

\section{Overplotting}

\begin{frame}{Problem and Three Solutions}
  \begin{center}
    \includegraphics[width=0.95\textwidth]{figures_w02m03/overplotting}
  \end{center}
\end{frame}

\begin{frame}[fragile]{Overplotting Solutions in R}
\begin{lstlisting}[style=rstyle]
# Fix 1: Reduce alpha (most common)
ggplot(df, aes(x, y)) +
  geom_point(alpha = 0.05, size = 0.5)     # near-transparent

# Fix 2: Jitter (for discrete x)
ggplot(df, aes(x = category, y = value)) +
  geom_jitter(width = 0.2, height = 0, alpha = 0.15, size = 0.8)

# Fix 3: Hexbin (for large n)
ggplot(df, aes(x, y)) +
  geom_hex(bins = 25) +
  scale_fill_viridis_c()                   # count encoded as colour

# Fix 4: 2D density contours
ggplot(df, aes(x, y)) +
  geom_density_2d_filled(alpha = 0.7) +
  geom_point(alpha = 0.02, size = 0.3)     # ghost points for context
\end{lstlisting}
  \begin{alertblock}{Pro-Tip}
    Rule of thumb: $n < 500$ use alpha; $500 < n < 5000$ use hex/density;
    $n > 5000$ use aggregation (\texttt{geom\_bin2d} or \texttt{geom\_hex}).
  \end{alertblock}
\end{frame}

\section{Special Geometries}

\begin{frame}{Annotation Geoms: text vs label}
  \begin{center}
    \includegraphics[width=0.92\textwidth]{figures_w02m03/annotation_geoms}
  \end{center}
  \begin{exampleblock}{Data Insight}
    \texttt{geom\_text()} renders plain text (no background). \texttt{geom\_label()}
    adds a filled box --- better visibility on cluttered backgrounds.
    For anti-overlap, use \texttt{ggrepel::geom\_text\_repel()} (R) or
    \texttt{adjustText} (Python).
  \end{exampleblock}
\end{frame}

\begin{frame}{Connector Geoms: Lollipop \& Dumbbell}
  \begin{center}
    \includegraphics[width=0.92\textwidth]{figures_w02m03/connector_geoms}
  \end{center}
\end{frame}

\begin{frame}[fragile]{Lollipop \& Dumbbell in R}
\begin{lstlisting}[style=rstyle]
# Lollipop: geom_segment + geom_point
ggplot(df, aes(x = fct_reorder(cat, val), y = val)) +
  geom_segment(aes(xend = cat, y = 0, yend = val), color = "grey70") +
  geom_point(size = 3, color = "#1565C0") +
  coord_flip() + theme_minimal()

# Dumbbell: two geom_point layers + geom_segment
ggplot(df, aes(y = fct_reorder(cat, val_2024))) +
  geom_segment(aes(x = val_2020, xend = val_2024,
                   yend = cat), color = "grey80", linewidth = 1.5) +
  geom_point(aes(x = val_2020), size = 3, color = "#1565C0") +
  geom_point(aes(x = val_2024), size = 3, color = "#E53935") +
  theme_minimal() +
  labs(x = "Value", y = NULL, title = "Change: 2020 vs 2024")
\end{lstlisting}
\end{frame}

\section{Synthesis}

\begin{frame}{Aesthetic Selection Checklist}
  \centering
  \begin{tabular}{lll}
    \toprule
    \textbf{Variable Type} & \textbf{Best Aesthetic} & \textbf{Fallback} \\
    \midrule
    Quantitative (primary) & x or y position & size \\
    Quantitative (secondary) & colour gradient & size \\
    Categorical ($\leq$ 6) & colour hue or shape & facet \\
    Categorical ($>$ 6) & facet\_wrap & colour + label \\
    Ordered categorical & colour luminance & size \\
    Binary & colour (2 hues) & shape or alpha \\
    \bottomrule
  \end{tabular}
  \vspace{6pt}
  \begin{alertblock}{Pro-Tip}
    Encode the most important variable on the highest-precision channel
    (position). Reserve colour for the second-most-important grouping.
  \end{alertblock}
\end{frame}

\begin{frame}{Key Takeaways}
  \begin{enumerate}
    \item Five core aesthetics: \textbf{x/y} (position), \textbf{color}
          (hue/gradient), \textbf{size} (area), \textbf{shape} (glyph,
          $\leq$6), \textbf{alpha} (transparency).
    \item Match the aesthetic to the variable type: continuous $\rightarrow$
          position or gradient; categorical $\rightarrow$ hue or shape.
    \item \textbf{Layering geoms} with \texttt{+} creates composite charts:
          point+smooth, boxplot+jitter, col+errorbar.
    \item \textbf{Overplotting} is solved by alpha ($n<500$), jitter
          (discrete x), hex/density ($n>500$).
    \item \textbf{Special geoms}: text/label for annotation,
          segment+point for lollipop and dumbbell charts.
  \end{enumerate}
\end{frame}

\begin{frame}{Essential Readings}
  \begin{itemize}
    \item Wickham, H. (2016). \emph{ggplot2}, Chapters 4--5
          (Mastering Aesthetics, Geom Catalogue).
    \item Wilke, C.~O. (2019). \emph{Fundamentals of Data Visualization},
          Chapters 2, 18 (Aesthetics, Handling Overlapping Points).
    \item Chang, W. (2024). \emph{R Graphics Cookbook}, 2nd ed., Chapters
          5--6 (Scatter Plots, Bar Charts).
  \end{itemize}
  \vspace{6pt}
  \textbf{Next Module}: Scales \& Coordinate Systems (W02-M04)
\end{frame}

\end{document}
